\chapter{Review of methods and tools}
\label{cha:methods-and-tools}

Two themes recur across the choices below. The first is a preference for offline, version-pinned artefacts over live services: a food knowledge base queried dozens of times while composing a single plan cannot depend on a rate-limited network call, and a result set that changes underneath a running experiment cannot be reproduced. The second is that a component's correctness guarantees matter more than its raw performance wherever a wrong answer would be indistinguishable from a right one, which is why a validated schema is preferred over a convention the caller is trusted to follow. Each section names the alternatives that were considered and the reason one of them was carried into \Cref{cha:design}. Concrete model identifiers appear in \Cref{cha:results}.

\section{Implementation platform}
\label{sec:tools-platform}

A dietary planning agent spends most of its wall-clock time waiting on network calls to language model providers and on local food knowledge base queries; it is not compute-bound. This makes the choice of ecosystem more important than raw execution speed. Python was chosen for its mature tooling for the two things this project spends the most engineering effort on: LLM tool-calling orchestration and structured output validation. Its ecosystem also supplies LLM client libraries for every major provider, an embedded analytical database with first-class bindings (\Cref{sec:tools-storage}), and a scientific-computing stack sufficient for the exact-arithmetic aggregation and plotting of the evaluation harness. A statically typed, compiled language would have offered better throughput and stronger compile-time guarantees, but a materially less mature ecosystem for orchestration and output validation.

\section{Agent framework}
\label{sec:tools-agent-framework}

Three broad approaches were available for orchestrating a language model that must call tools and produce a validated, typed output.

\begin{itemize}
    \item \textbf{Raw provider SDKs.} Offer maximum control and no framework dependency, but require re-implementing the same conversation-loop, retry and validation logic for every provider, since each provider's tool-calling wire format differs.
    \item \textbf{General-purpose orchestration frameworks (e.g.\ LangChain,\footnote{\url{https://python.langchain.com/} (accessed 2026-09-13).} LangGraph\footnote{\url{https://langchain-ai.github.io/langgraph/} (accessed 2026-09-13).}).} Abstract over multiple providers and add constructs for multi-step or graph-structured workflows, at the cost of a large surface area aimed at use cases (long-running stateful graphs, complex branching) that this system's mostly linear pipeline (\Cref{sec:design-pipeline}) does not need. Their structured-output support has also historically been an add-on to a chat-completion abstraction.
    \item \textbf{A schema-first, type-centric agent library (pydantic-ai).\footnote{\url{https://ai.pydantic.dev/} (accessed 2026-09-13).}} Treats a strongly typed data-validation model as the primary interface for both structured output and tool signatures, delegates provider-specific formatting to per-provider adapters, and exposes an output-validator hook that can reject a structurally valid but semantically wrong response and trigger an automatic retry with an explanatory message.
\end{itemize}

pydantic-ai was chosen. Its output-validator hook is the mechanism the no-invented-foods contract in \Cref{sec:design-no-invented-foods} depends on directly: every plan is validated against the real food knowledge base before acceptance, and a failure produces a targeted retry, not a silent pass-through. Its provider adapters also make \ref{req:provider-agnostic}~(provider-agnostic model access) close to free, since switching the underlying model is a configuration change.

\section{Model access}
\label{sec:tools-model-access}

Three ways of obtaining model inference were considered. Direct per-provider API access offers the lowest latency and tightest control over provider-specific features, but ties the codebase to one vendor's request format per model family, conflicting with \ref{req:provider-agnostic}. Self-hosted inference removes any external dependency and per-token cost, but requires provisioning GPU infrastructure sized for a capable enough model, an investment orthogonal to this thesis's research question. A routing service exposing a single, provider-agnostic API over many hosted models was chosen instead: it lets the model comparison in \Cref{cha:results} span multiple providers without separate integration code, and lets the ablation grid be re-run against a different model by changing a configuration value.

Contemporary reasoning models expose a second, independent control: a \emph{reasoning-effort} setting governing how much internal deliberation the provider allows before committing to an answer, typically a small ordinal scale from minimal to very high. It trades latency for deliberation within one model, and \ref{req:provider-agnostic}~requires it to be configurable. It is therefore treated as part of the model identifier, so a model and an effort level travel together as one configuration value, and it becomes the second axis of the capability comparison in \Cref{sec:methodology-model-grid}.

\section{Food composition data access}
\label{sec:tools-food-data}

A grounding strategy of the kind described in \Cref{sec:background-tool-use,sec:background-self-correction} is only as good as the knowledge base it grounds against. Two large, openly available sources cover complementary needs. The USDA FoodData Central Foundation Foods and SR Legacy datasets\footnote{\url{https://fdc.nal.usda.gov/} (accessed 2026-09-13).} provide laboratory-analysed nutrient profiles for generic, minimally processed foods, with broad applicability but no coverage of branded products. Open Food Facts\footnote{\url{https://world.openfoodfacts.org/} (accessed 2026-09-13).} is a crowd-sourced, international database of branded, packaged products; it has far denser coverage of packaged food, but far less consistent completeness of any individual nutrient field, since it depends on what each contributor transcribed. How the two are combined, and how an absent field is kept from silently becoming a zero, is taken up in \Cref{sec:design-food-db,sec:impl-missing-nutrients}.

Each source can be accessed through a live query API or a bulk offline export. A live API always reflects the provider's latest data, but introduces a network dependency and rate limits into every food lookup, which a system that may issue dozens of lookups per plan cannot tolerate on latency or reliability grounds. A commercial nutrition data API was also considered, but rejected on licensing and reproducibility grounds: a bulk offline export can be version-pinned and redistributed alongside the results it produced, which a commercial API's licence terms and rate limits make difficult. Bulk offline exports of both sources were chosen, downloaded once during a setup step (\Cref{sec:design-food-db}) and never queried live at request time.

\section{Storage and indexing engine}
\label{sec:tools-storage}

The food knowledge base must support three access patterns: exact and fuzzy text search over food names, similarity search over a per-food embedding vector, and relational lookups by identifier. Three engines were compared.

\begin{itemize}
    \item \textbf{SQLite with the FTS5 extension.\footnote{\url{https://www.sqlite.org/fts5.html} (accessed 2026-09-13).}} Zero-configuration and sufficient for lexical search, but with no first-party vector index; nearest-neighbour search would need a third-party extension, versioned and tested outside the engine itself.
    \item \textbf{PostgreSQL with the pgvector extension.\footnote{\url{https://github.com/pgvector/pgvector} (accessed 2026-09-13).}} A mature relational database with full-text and vector search as extensions, and the natural choice under concurrent multi-client write traffic -- but a server process to provision and keep running, disproportionate for a dataset built once during setup and only ever read.
    \item \textbf{DuckDB~\cite{raasveldt2019duckdb}.} An embedded analytical database that requires no server process, ships with a Python binding, and is optimised for read-mostly analytical queries. Full-text search and vector similarity are later first-party extensions~\cite{duckdb2026fts,duckdb2026vss}, not part of the 2019 system the paper introduces.
\end{itemize}

DuckDB was chosen. Its embedded model matches a food knowledge base built once offline and shipped read-only, and the two extensions cover both retrieval channels of \Cref{sec:tools-retrieval} without a second engine.

\section{Hybrid retrieval strategy}
\label{sec:tools-retrieval}

Searching the food knowledge base for, say, \enquote{mozzarella} must succeed whether the query is an exact lexical match for a product name or a semantic near-match for a differently worded description. Two retrieval families address these cases with complementary strengths.

Lexical ranking (the BM25 scoring function~\cite{robertson2009probabilistic}) is fast and needs no learned model, but is blind to synonymy: a query for \enquote{courgette} will not match a product labelled only \enquote{zucchini}. Dense retrieval, embedding query and catalogue entries into one vector space and ranking by similarity, captures such semantic relationships, but can under-rank a close lexical match that embeds slightly off-distribution, and is more expensive to compute and index.

Combining the two rankings requires a fusion strategy. Score-based fusion (a weighted sum of raw similarity scores) requires normalising the scores onto a comparable scale, which is fragile because BM25 and cosine-similarity scores have different, corpus-dependent distributions. Reciprocal rank fusion sidesteps this by combining only each result's \emph{rank} within its own list, via $\mathrm{RRF}(d) = \sum_i \frac{1}{k + \mathrm{rank}_i(d)}$ summed over each list $i$ a document $d$ appears in, with $k$ a small constant (conventionally \num{60}) that dampens the impact of a single high ranking. Without training data or score calibration it exceeded Condorcet Fuse on every reported collection and CombMNZ on all but two: CombMNZ was ahead by \num{0.003} mean average precision on TREC~3, for which no per-collection significance test is reported, and by a margin the authors put at $p \approx \num{0.2}$ on LETOR~3. It also outperformed the best individual ranking in each experiment except TREC~9, where that ranking was produced with a human in the loop~\cite{cormack2009reciprocal}. It was chosen for exactly this rank-only, calibration-free property.

\Cref{fig:diagram-rrf-example} works the formula through on one query. Beyond calibration-freedom, it exposes the decisive behaviour: a candidate both channels rank highly is more likely the intended food than one a single channel ranks first, and rank-only fusion prefers it without being told to. \Cref{sec:design-hybrid-search} describes the scheme as wired into the lookup tool.

\begin{figure}[htbp]
    \centering
    \resizebox{0.92\textwidth}{!}{\begin{tikzpicture}[
    list/.style={draw, rounded corners, align=left, inner sep=2mm, font=\small},
    box/.style={draw, rounded corners, align=center, inner sep=2mm, font=\small},
    title/.style={font=\small\bfseries, align=center},
    arr/.style={-{Latex[length=2mm]}},
]

\node[list, anchor=north west] (lex) at (0mm, 0mm) {
    \begin{tabular}{@{}r@{~~}l@{~~}r@{}}
        \toprule
        rank & catalogue entry & \textcolor{black!40}{\makecell{BM25\\ score}}\\
        \midrule
        1 & mozzarella, \qty{125}{\gram}  & \textcolor{black!40}{9.1}\\
        2 & buffalo mozzarella, fresh     & \textcolor{black!40}{7.4}\\
        3 & mozzarella pizza topping      & \textcolor{black!40}{6.2}\\
        4 & low-moisture mozzarella       & \textcolor{black!40}{5.5}\\
        \bottomrule
    \end{tabular}
};

\node[list, anchor=north west] (dense) at ($(lex.north east) + (14mm, 0mm)$) {
    \begin{tabular}{@{}r@{~~}l@{~~}r@{}}
        \toprule
        rank & catalogue entry & \textcolor{black!40}{\makecell{cosine\\ distance}}\\
        \midrule
        1 & fresh cheese in brine         & \textcolor{black!40}{0.18}\\
        2 & buffalo mozzarella, fresh     & \textcolor{black!40}{0.22}\\
        3 & low-moisture mozzarella       & \textcolor{black!40}{0.29}\\
        4 & mozzarella, \qty{125}{\gram}  & \textcolor{black!40}{0.34}\\
        \bottomrule
    \end{tabular}
};

\node[title, anchor=south] at ($(lex.north) + (0mm, 1.5mm)$) {Lexical channel};
\node[title, anchor=south] at ($(dense.north) + (0mm, 1.5mm)$) {Dense channel};
\node[font=\small, anchor=south] at ($(lex.north)!0.5!(dense.north) + (0mm, 9mm)$)
    {query: \texttt{mozzarella}};

\node[box, anchor=north] (rrf) at ($(lex.south)!0.5!(dense.south) + (0mm, -14mm)$)
    {$\sum_i \frac{1}{k + \mathrm{rank}_i(d)}$, $k = 60$};
\node[title, anchor=south] at ($(rrf.north) + (0mm, 1.5mm)$) {Reciprocal rank fusion};

\draw[arr] (lex.south) |- (rrf.west);
\draw[arr] (dense.south) |- (rrf.east);

\node[list, anchor=north] (fused) at ($(rrf.south) + (0mm, -10mm)$) {
    \begin{tabular}{@{}r@{~~}l@{~~}c@{~~}r@{}}
        \toprule
        rank & catalogue entry & \makecell{ranks\\ (lex., dense)} & \makecell{fused\\ score}\\
        \midrule
        1 & buffalo mozzarella, fresh    & 2,\,2   & 0.0323\\
        2 & mozzarella, \qty{125}{\gram} & 1,\,4   & 0.0320\\
        3 & low-moisture mozzarella      & 4,\,3   & 0.0315\\
        4 & fresh cheese in brine        & --,\,1  & 0.0164\\
        5 & mozzarella pizza topping     & 3,\,--  & 0.0159\\
        \bottomrule
    \end{tabular}
};
\node[title, anchor=south] at ($(fused.north) + (0mm, 1.5mm)$) {Fused ranking};

\draw[arr] (rrf.south) -- (fused.north);

\end{tikzpicture}
}
    \caption{Reciprocal rank fusion on a mozzarella query.}
    \label{fig:diagram-rrf-example}
\end{figure}

\section{Embedding model and vector index}
\label{sec:tools-embeddings}

Dense retrieval requires an embedding model to turn a food description into a vector, and an index to search that vector space at scale. For the index, Hierarchical Navigable Small World graphs (HNSW) were chosen over a brute-force scan: under the paper's bounded-degree assumptions the query cost grows logarithmically, not linearly, with the size of the knowledge base, at a small, tunable cost in recall~\cite{malkov2018efficient}. Those assumptions were demonstrated on low-dimensional data; this catalogue uses 1024-dimensional embeddings, so the practical gain is sub-linear search with a native index type in the storage engine of \Cref{sec:tools-storage}, avoiding a second vector database.

For the embedding model, three options existed: a hosted embedding API, a small English-only local model, and a larger multilingual local model. A hosted API was rejected for the same latency and reproducibility reasons as a live food-data API in \Cref{sec:tools-food-data}. A small English-only model would be cheapest to run but would degrade on the many non-English names in an internationally sourced catalogue. BGE-M3, a multilingual model capable of dense, sparse and multi-vector retrieval from one set of weights, trained to generalise across more than a hundred languages and document lengths~\cite{chen2024bgem3}, was chosen instead, run locally so that embedding the catalogue during setup and embedding a query at request time both stay off the network.

\section{Self-correction scheme}
\label{sec:tools-self-correction}

\Cref{sec:background-self-correction} identified two points on a spectrum of self-correction: a critique-and-refine loop grounded in a model's own re-reading of its output, and one grounded in feedback from an external tool. A fully tool-grounded critic, in the style of \textsc{CRITIC}~\cite{gou2024critic}, whose review is checked against external tool feedback, was considered but set aside for a narrower design: the critic is tool-less but is still shown the deterministic totaller's output when the totaller is enabled. This keeps most of the grounding benefit while giving the critic a fixed, predictable set of inputs and not the ability to trigger unbounded retrieval. \Cref{sec:design-reflection} describes the resulting loop, and \Cref{cha:results} reports how much of the grounding gain survives when the critic, not only the generator, sees the computed totals.

\Cref{fig:diagram-self-correction} traces the reference design beside the three ablation variants of this thesis, one of which has no review loop.

\begin{figure}[htbp]
    \centering
    \resizebox{\textwidth}{!}{\begin{tikzpicture}[
    x=1mm, y=1mm,
    toollane/.style={draw, rounded corners=1mm, anchor=west, text width=110mm,
                     minimum height=5mm, inner sep=1mm, font=\small, align=center},
    toolcall/.style={{Latex[length=2mm]}-{Latex[length=2mm]}, draw=black!55, dashed},
    toolcallensured/.style={{Latex[length=2mm]}-{Latex[length=2mm]}},
    box/.style={draw, rounded corners, anchor=west, text width=24mm,
                minimum height=11mm, inner sep=1.2mm, font=\small, align=center},
    edgelab/.style={font=\small, text=black!60},
    note/.style={font=\small, text=black!60, align=center, text width=80mm},
    title/.style={font=\small\bfseries, anchor=south west},
    frame/.style={draw=black!35, rounded corners},
    arr/.style={-{Latex[length=2mm]}},
]

% Shared local coordinates in every panel:
%   tool bar at y = 0 spanning x = 0..112
%   prompt enters from x = -22; output leaves the frame at x = 116
%   Agent (0, -20), Critic (42, -20), Revise / Refiner (84, -20)
%   boxes are 26.4 wide, so the 15.6 gaps hold the edge labels

% ===== panel 1: CRITIC =====
\begin{scope}[yshift=0mm]
    \draw[frame] (-4, 6) rectangle (116, -42);
    \node[title] at (-4, 8) {CRITIC, the reference design};

    \node[toollane] at (0, 0) {External verification tool};
    \node[box] (c_agent) at (0, -20) {Agent};
    \node[box] (c_critic) at (42, -20) {Critic};
    \node[box] (c_revise) at (84, -20) {Revise};
    \draw[arr] (-22, -20) -- node[edgelab, above] {prompt} (c_agent.west);
    \draw[arr] (c_agent.east) -- node[edgelab, above] {draft} (c_critic.west);
    \draw[arr] (c_critic.east) -- node[edgelab, above] {feedback} (c_revise.west);
    \draw[toolcallensured] (55, -2.5) -- node[edgelab, right, pos=0.45] {ensured} (55, -14.5);

    \draw[arr] (c_revise.south) -- ++(0, -5) -| ([xshift=8mm]c_critic.south);
    \node[edgelab, anchor=north] at (75.6, -31) {revised result};

    \draw[arr] (c_critic.south) -- ++(0, -11) -- node[edgelab, below, pos=0.88] {result} (138, -36.5);
\end{scope}

% ===== panel 2: totaller-only =====
\begin{scope}[yshift=-62mm]
    \draw[frame] (-4, 6) rectangle (116, -29);
    \node[title] at (-4, 8) {\texttt{totaller-only}};

    \node[toollane] at (0, 0) {Totaller};
    \node[box] (t_agent) at (0, -20) {Nutrition Agent};
    \draw[arr] (-22, -20) -- node[edgelab, above] {prompt} (t_agent.west);
    \draw[arr] (t_agent.east) -- node[edgelab, above, pos=0.88] {plan} (138, -20);
    \draw[toolcall] (6, -14.5) -- (6, -2.5);
    \draw[toolcall] (13, -14.5) -- (13, -2.5);
    \draw[toolcall] (21, -14.5) -- (21, -2.5);
\end{scope}

% ===== panel 3: reflection-only =====
\begin{scope}[yshift=-110mm]
    \draw[frame] (-4, 6) rectangle (116, -42);
    \node[title] at (-4, 8) {\texttt{reflection-only}};

    \node[note] at (56, 0) {no computed totals available to any role};

    \node[box] (r_agent) at (0, -20) {Nutrition Agent};
    \node[box] (r_critic) at (42, -20) {Critic};
    \node[box] (r_refiner) at (84, -20) {Refiner};
    \draw[arr] (-22, -20) -- node[edgelab, above] {prompt} (r_agent.west);
    \draw[arr] (r_agent.east) -- node[edgelab, above] {draft plan} (r_critic.west);
    \draw[arr] (r_critic.east) -- node[edgelab, above] {feedback} (r_refiner.west);

    \draw[arr] (r_refiner.south) -- ++(0, -5) -| ([xshift=8mm]r_critic.south);
    \node[edgelab, anchor=north] at (75.6, -31) {refined plan};

    \draw[arr] (r_critic.south) -- ++(0, -11) -- node[edgelab, below, pos=0.88] {plan} (138, -36.5);
\end{scope}

% ===== panel 4: full =====
\begin{scope}[yshift=-171mm]
    \draw[frame] (-4, 6) rectangle (116, -42);
    \node[title] at (-4, 8) {\texttt{full}};

    \node[toollane] at (0, 0) {Totaller};
    \node[box] (f_agent) at (0, -20) {Nutrition Agent};
    \node[box] (f_critic) at (42, -20) {Critic};
    \node[box] (f_refiner) at (84, -20) {Refiner};
    \draw[arr] (-22, -20) -- node[edgelab, above] {prompt} (f_agent.west);
    \draw[arr] (f_agent.east) -- node[edgelab, above] {draft plan} (f_critic.west);
    \draw[arr] (f_critic.east) -- node[edgelab, above] {feedback} (f_refiner.west);

    \draw[toolcall] (6, -14.5) -- (6, -2.5);
    \draw[toolcall] (13, -14.5) -- (13, -2.5);
    \draw[toolcall] (21, -14.5) -- (21, -2.5);
    \draw[toolcall] (104, -14.5) -- (104, -2.5);

    \draw[toolcallensured] (55, -2.5) -- node[edgelab, right, pos=0.45] {ensured} (55, -14.5);
    \draw[toolcallensured] (90, -2.5) -- node[edgelab, left, pos=0.45] {ensured} (90, -14.5);

    \draw[arr] (f_refiner.south) -- ++(0, -5) -| ([xshift=8mm]f_critic.south);
    \node[edgelab, anchor=north] at (75.6, -31) {refined plan};

    \draw[arr] (f_critic.south) -- ++(0, -11) -- node[edgelab, below, pos=0.88] {plan} (138, -36.5);
\end{scope}

\end{tikzpicture}
}
    \caption{Critique-and-refine configurations used in this thesis compared with CRITIC design.}
    \label{fig:diagram-self-correction}
\end{figure}

\section{Evaluation tooling}
\label{sec:tools-evaluation}

The evaluation harness needs three capabilities distinct from the runtime it evaluates: driving a large, resumable grid of pipeline runs across models, configurations and scenarios; scoring both the numeric and open-ended, natural-language aspects of a run; and turning the scores into readable summaries and figures. It was built as a separate command-line application, following the package-layering rule that keeps evaluation-only code out of the runtime path (\Cref{sec:design-layering}). The qualitative scoring stage reuses the schema-first agent framework of \Cref{sec:tools-agent-framework}, configured with a per-criterion rubric and following the G-Eval pattern of \Cref{sec:background-evaluation}. The judge protocol is specified in \Cref{sec:methodology-judge-protocol}.

\section{Summary of the choices made}
\label{sec:tools-summary}

The decisions of this chapter are the contract the design in \Cref{cha:design} is written against. Each row of \Cref{tab:tool-choices} names the component that later chapters treat as given, and the alternative it was chosen over, so that a later change of tool can be traced to a rejected option and not to an unspoken default. Two preferences run through the table. Offline, version-pinned artefacts were preferred wherever a live service would make a grid irreproducible. Correctness guarantees were preferred wherever a wrong answer would be indistinguishable from a right one. Those preferences, and not raw throughput, are why the chapters that follow assume an embedded food knowledge base, offline bulk data, and a schema-first agent framework.

\begin{table}[htbp]
    \centering
    \small
    \begin{tabularx}{\textwidth}{@{}
    >{\raggedright\arraybackslash}p{0.24\textwidth}
    >{\raggedright\arraybackslash}p{0.28\textwidth}
    >{\raggedright\arraybackslash}X
    @{}}
    \toprule
    Name and section & Chosen & Rejected alternatives \\
    \midrule
    Implementation platform\newline
    \textit{\footnotesize\Cref{sec:tools-platform}}
        & Python
        & Statically typed, compiled language \\
    \addlinespace
    Agent framework\newline
    \textit{\footnotesize\Cref{sec:tools-agent-framework}}
        & pydantic-ai
        & Raw provider SDKs\newline
          General-purpose orchestration framework \\
    \addlinespace
    Model access\newline
    \textit{\footnotesize\Cref{sec:tools-model-access}}
        & Routing service over many providers
        & Direct per-provider API\newline
          Self-hosted open-weight inference \\
    \addlinespace
    Food data access\newline
    \textit{\footnotesize\Cref{sec:tools-food-data}}
        & Bulk, version-pinned offline exports
        & Live query APIs\newline
          Commercial nutrition data API \\
    \addlinespace
    Storage engine\newline
    \textit{\footnotesize\Cref{sec:tools-storage}}
        & DuckDB, embedded, full-text and vector
        & SQLite with FTS5\newline
          PostgreSQL with pgvector \\
    \addlinespace
    Retrieval channel\newline
    \textit{\footnotesize\Cref{sec:tools-retrieval}}
        & Both channels, in parallel
        & Lexical ranking alone (BM25)\newline
          Dense ranking alone \\
    \addlinespace
    Fusion strategy\newline
    \textit{\footnotesize\Cref{sec:tools-retrieval}}
        & Reciprocal rank fusion
        & Weighted sum of raw scores \\
    \addlinespace
    Embedding model\newline
    \textit{\footnotesize\Cref{sec:tools-embeddings}}
        & BGE-M3, multilingual, local
        & Hosted embedding API\newline
          Small, English-only local model \\
    \addlinespace
    Vector index\newline
    \textit{\footnotesize\Cref{sec:tools-embeddings}}
        & HNSW proximity graph
        & Exact brute-force scan \\
    \addlinespace
    Self-correction\newline
    \textit{\footnotesize\Cref{sec:tools-self-correction}}
        & Tool-less critic shown computed totals
        & Critic re-reading its own output only\newline
          Critic issuing its own tool calls \\
    \addlinespace
    Judge protocol\newline
    \textit{\footnotesize\Cref{sec:tools-evaluation}}
        & Two vendors' judges, repeated scoring
        & One judge, one call per plan \\
    \addlinespace
    Evaluation harness\newline
    \textit{\footnotesize\Cref{sec:tools-evaluation}}
        & Separate command-line application
        & Evaluation code inside the runtime \\
    \bottomrule
\end{tabularx}

    \caption{Technology decisions carried into \Cref{cha:design}.}
    \label{tab:tool-choices}
\end{table}
